\documentclass[10pt,a4paper]{amsart}
\usepackage{amsmath,amssymb,amsthm,graphicx}
% Retained preamble for the Stacks Project chapter snapshot in this directory.
% The chapter's own preamble.tex is a copy of the shared Stacks Project preamble; the block
% below is sliced MECHANICALLY, line for line, from it. Nothing was added, removed or
% redefined; the wrapper class, the xr-hyper/hyperref wiring and the externaldocument list
% are omitted because the assembler supplies the document class.
%
% lines taken from preamble.tex: 3, 5, 9, 11, 12, 156, 157, 158, ... (51 lines total)
%
% package coverage restored from the chapter's own preamble (2026-10-13): xy, xy-2cell

\usepackage[all]{xy}
\xyoption{2cell}
\UseAllTwocells

\newenvironment{reference}{\comment}{\endcomment}
\newenvironment{slogan}{\comment}{\endcomment}
\newenvironment{history}{\comment}{\endcomment}
\theoremstyle{plain}
\newtheorem{theorem}[subsection]{Theorem}
\newtheorem{proposition}[subsection]{Proposition}
\newtheorem{lemma}[subsection]{Lemma}
\theoremstyle{definition}
\newtheorem{definition}[subsection]{Definition}
\newtheorem{example}[subsection]{Example}
\newtheorem{exercise}[subsection]{Exercise}
\newtheorem{situation}[subsection]{Situation}
\theoremstyle{remark}
\newtheorem{remark}[subsection]{Remark}
\newtheorem{remarks}[subsection]{Remarks}
\numberwithin{equation}{subsection}
\def\lim{\mathop{\mathrm{lim}}\nolimits}
\def\colim{\mathop{\mathrm{colim}}\nolimits}
\def\Spec{\mathop{\mathrm{Spec}}}
\def\Hom{\mathop{\mathrm{Hom}}\nolimits}
\def\Ext{\mathop{\mathrm{Ext}}\nolimits}
\def\SheafHom{\mathop{\mathcal{H}\!\mathit{om}}\nolimits}
\def\SheafExt{\mathop{\mathcal{E}\!\mathit{xt}}\nolimits}
\def\Sch{\mathit{Sch}}
\def\Mor{\mathop{\mathrm{Mor}}\nolimits}
\def\Ob{\mathop{\mathrm{Ob}}\nolimits}
\def\Sh{\mathop{\mathit{Sh}}\nolimits}
\def\NL{\mathop{N\!L}\nolimits}
\def\CH{\mathop{\mathrm{CH}}\nolimits}
\def\proetale{{pro\text{-}\acute{e}tale}}
\def\etale{{\acute{e}tale}}
\def\QCoh{\mathit{QCoh}}
\def\Ker{\mathop{\mathrm{Ker}}}
\def\Im{\mathop{\mathrm{Im}}}
\def\Coker{\mathop{\mathrm{Coker}}}
\def\Coim{\mathop{\mathrm{Coim}}}
\DeclareMathSymbol{\boxtimes}{\mathbin}{AMSa}{"02}
\def\QCohstack{\mathcal{QC}\!\mathit{oh}}
\def\Cohstack{\mathcal{C}\!\mathit{oh}}
\def\Spacesstack{\mathcal{S}\!\mathit{paces}}
\def\Quotfunctor{\mathrm{Quot}}
\def\Hilbfunctor{\mathrm{Hilb}}
\def\Curvesstack{\mathcal{C}\!\mathit{urves}}
\def\Polarizedstack{\mathcal{P}\!\mathit{olarized}}
\def\Complexesstack{\mathcal{C}\!\mathit{omplexes}}
\def\Pic{\mathop{\mathrm{Pic}}\nolimits}
\def\Picardstack{\mathcal{P}\!\mathit{ic}}
\def\Picardfunctor{\mathrm{Pic}}
\def\Deformationcategory{\mathcal{D}\!\mathit{ef}}
\usepackage{cleveref}
\crefname{originallemma}{original lemma}{original lemmas}
\Crefname{originallemma}{Original lemma}{Original lemmas}
\makeatletter
\newlabel{algebraic-lemma-2-fibre-product-general}{{04TF}{source}{}{sourcepointer-1}{}}
\newlabel{algebraic-lemma-2-fibre-product-general@cref}{{[originallemma][0][]04TF}{[1][0][]source}{}{}{}}
\makeatother
\crefname{originallemma}{original lemma}{original lemmas}
\Crefname{originallemma}{Original lemma}{Original lemmas}
\makeatletter
\newlabel{algebraic-lemma-stack-presentation}{{04T5}{source}{}{sourcepointer-2}{}}
\newlabel{algebraic-lemma-stack-presentation@cref}{{[originallemma][0][]04T5}{[1][0][]source}{}{}{}}
\makeatother
\crefname{originallemma}{original lemma}{original lemmas}
\Crefname{originallemma}{Original lemma}{Original lemmas}
\makeatletter
\newlabel{equation-functoriality}{{07WJ}{source}{}{sourcepointer-3}{}}
\newlabel{equation-functoriality@cref}{{[originallemma][0][]07WJ}{[1][0][]source}{}{}{}}
\makeatother
\crefname{originallemma}{original lemma}{original lemmas}
\Crefname{originallemma}{Original lemma}{Original lemmas}
\makeatletter
\newlabel{formal-defos-lemma-deformation-functor-diagonal}{{06L6}{source}{}{sourcepointer-4}{}}
\newlabel{formal-defos-lemma-deformation-functor-diagonal@cref}{{[originallemma][0][]06L6}{[1][0][]source}{}{}{}}
\makeatother
\crefname{originallemma}{original lemma}{original lemmas}
\Crefname{originallemma}{Original lemma}{Original lemmas}
\makeatletter
\newlabel{formal-defos-lemma-inf-aut-change-of-field}{{07WC}{source}{}{sourcepointer-5}{}}
\newlabel{formal-defos-lemma-inf-aut-change-of-field@cref}{{[originallemma][0][]07WC}{[1][0][]source}{}{}{}}
\makeatother
\crefname{originallemma}{original lemma}{original lemmas}
\Crefname{originallemma}{Original lemma}{Original lemmas}
\makeatletter
\newlabel{formal-defos-lemma-k-linear-differential}{{06IJ}{source}{}{sourcepointer-6}{}}
\newlabel{formal-defos-lemma-k-linear-differential@cref}{{[originallemma][0][]06IJ}{[1][0][]source}{}{}{}}
\makeatother
\crefname{originallemma}{original lemma}{original lemmas}
\Crefname{originallemma}{Original lemma}{Original lemmas}
\makeatletter
\newlabel{formal-defos-lemma-smooth-morphism-essentially-surjective}{{06HN}{source}{}{sourcepointer-7}{}}
\newlabel{formal-defos-lemma-smooth-morphism-essentially-surjective@cref}{{[originallemma][0][]06HN}{[1][0][]source}{}{}{}}
\makeatother
\crefname{originallemma}{original lemma}{original lemmas}
\Crefname{originallemma}{Original lemma}{Original lemmas}
\makeatletter
\newlabel{formal-defos-lemma-tangent-space-change-of-field}{{07WB}{source}{}{sourcepointer-8}{}}
\newlabel{formal-defos-lemma-tangent-space-change-of-field@cref}{{[originallemma][0][]07WB}{[1][0][]source}{}{}{}}
\makeatother
\crefname{originallemma}{original lemma}{original lemmas}
\Crefname{originallemma}{Original lemma}{Original lemmas}
\makeatletter
\newlabel{lemma-algebraic-stack-RS}{{07WQ}{source}{}{sourcepointer-9}{}}
\newlabel{lemma-algebraic-stack-RS@cref}{{[originallemma][0][]07WQ}{[1][0][]source}{}{}{}}
\makeatother
\crefname{originallemma}{original lemma}{original lemmas}
\Crefname{originallemma}{Original lemma}{Original lemmas}
\makeatletter
\newlabel{lemma-change-of-field}{{07WX}{source}{}{sourcepointer-10}{}}
\newlabel{lemma-change-of-field@cref}{{[originallemma][0][]07WX}{[1][0][]source}{}{}{}}
\makeatother
\crefname{originallemma}{original lemma}{original lemmas}
\Crefname{originallemma}{Original lemma}{Original lemmas}
\makeatletter
\newlabel{lemma-deformation-category}{{07WU}{source}{}{sourcepointer-11}{}}
\newlabel{lemma-deformation-category@cref}{{[originallemma][0][]07WU}{[1][0][]source}{}{}{}}
\makeatother
\crefname{originallemma}{original lemma}{original lemmas}
\Crefname{originallemma}{Original lemma}{Original lemmas}
\makeatletter
\newlabel{lemma-fibre-product-deformation-categories}{{07WL}{source}{}{sourcepointer-12}{}}
\newlabel{lemma-fibre-product-deformation-categories@cref}{{[originallemma][0][]07WL}{[1][0][]source}{}{}{}}
\makeatother
\crefname{originallemma}{original lemma}{original lemmas}
\Crefname{originallemma}{Original lemma}{Original lemmas}
\makeatletter
\newlabel{lemma-formally-smooth-on-deformation-categories}{{07WK}{source}{}{sourcepointer-13}{}}
\newlabel{lemma-formally-smooth-on-deformation-categories@cref}{{[originallemma][0][]07WK}{[1][0][]source}{}{}{}}
\makeatother
\crefname{originallemma}{original lemma}{original lemmas}
\Crefname{originallemma}{Original lemma}{Original lemmas}
\makeatletter
\newlabel{remark-deformation-category-implies}{{07WV}{source}{}{sourcepointer-14}{}}
\newlabel{remark-deformation-category-implies@cref}{{[originallemma][0][]07WV}{[1][0][]source}{}{}{}}
\makeatother
\crefname{originallemma}{original lemma}{original lemmas}
\Crefname{originallemma}{Original lemma}{Original lemmas}
\makeatletter
\newlabel{section-predeformation-categories}{{07T2}{source}{}{sourcepointer-15}{}}
\newlabel{section-predeformation-categories@cref}{{[originallemma][0][]07T2}{[1][0][]source}{}{}{}}
\makeatother
\crefname{originallemma}{original lemma}{original lemmas}
\Crefname{originallemma}{Original lemma}{Original lemmas}
\makeatletter
\newlabel{spaces-over-fields-lemma-smooth-separable-closed-points-dense}{{07W4}{source}{}{sourcepointer-16}{}}
\newlabel{spaces-over-fields-lemma-smooth-separable-closed-points-dense@cref}{{[originallemma][0][]07W4}{[1][0][]source}{}{}{}}
\makeatother
\title{Finite-dimensionality of tangent and infinitesimal automorphism spaces}
\author{The Stacks Project Authors}
\date{stacks-project:artin.tex:local-snapshot-20261012; GNU FDL 1.2 or later}
\begin{document}
\maketitle
\noindent\textit{Editorial note.} Statement span is lines 1029-1042 of sources/stacks-artin/pinned-original-source.tex; proof span is lines 1044-1151, the complete author proof (106 body lines).\par\medskip
\noindent\textit{Editorial note.} Navigation stubs: 16 cited-but-not-included labels are supplied as non-load-bearing pointers with their pinned tags where the snapshot tag index has them.\par\medskip
\noindent\textit{Editorial note.} extra\_packages is empty on purpose: the retained preamble already loads xy with options and a second declaration clashes in the assembler wrapper.\par\medskip
\noindent\textit{Editorial note.} Text is retained verbatim, including the author's own typos and placeholders.\par\medskip
\noindent\textit{Editorial note.} Queue index 5 of queue\_artin.json (pb=4256, stmt 1029-1042, proof 1044-1151).\par\medskip
\noindent\textit{Editorial note.} stacks-artin wave ar01 (2026-10-14): fresh adjudication accept per lane-s3/artin\_judgments.md; pool per Chief 10:5x direct dispatch (idle-first over S12's 01:5x note, SMC precedent); artin.tex 99 rows/67 with proof judged 46 accept / 21 reject / 5 AGP-overlap avoidance rows; patch-sd-envs contract spans scanned. Named omission: 'Some details omitted' at the final dimension count - routine bookkeeping.\par\medskip
\noindent\textit{Allowed standard prerequisites (editorial).}\begin{itemize}
\item Explicit prerequisite boundary. The cited inputs are prerequisites; the argument in the retained span is not replaced by any cited result.
\end{itemize}\medskip
\ifdefined\hypertarget\hypertarget{sourcepointer-1}{}\fi
\noindent\textit{Source reference.} Original lemma 04TF: sources/stacks-artin/pinned-original-source.tex, source lines 1151--1151. Cited by the retained statement or proof but not part of the collected spans; supplied only as a navigation pointer. The argument does not depend on its content.\par\medskip
\ifdefined\hypertarget\hypertarget{sourcepointer-2}{}\fi
\noindent\textit{Source reference.} Original lemma 04T5: sources/stacks-artin/pinned-original-source.tex, source lines 1151--1151. Cited by the retained statement or proof but not part of the collected spans; supplied only as a navigation pointer. The argument does not depend on its content.\par\medskip
\ifdefined\hypertarget\hypertarget{sourcepointer-3}{}\fi
\noindent\textit{Source reference.} Original lemma 07WJ: sources/stacks-artin/pinned-original-source.tex, source lines 528--533. Cited by the retained statement or proof but not part of the collected spans; supplied only as a navigation pointer. The argument does not depend on its content.\par\medskip
\ifdefined\hypertarget\hypertarget{sourcepointer-4}{}\fi
\noindent\textit{Source reference.} Original lemma 06L6: sources/stacks-artin/pinned-original-source.tex, source lines 1151--1151. Cited by the retained statement or proof but not part of the collected spans; supplied only as a navigation pointer. The argument does not depend on its content.\par\medskip
\ifdefined\hypertarget\hypertarget{sourcepointer-5}{}\fi
\noindent\textit{Source reference.} Original lemma 07WC: sources/stacks-artin/pinned-original-source.tex, source lines 1151--1151. Cited by the retained statement or proof but not part of the collected spans; supplied only as a navigation pointer. The argument does not depend on its content.\par\medskip
\ifdefined\hypertarget\hypertarget{sourcepointer-6}{}\fi
\noindent\textit{Source reference.} Original lemma 06IJ: sources/stacks-artin/pinned-original-source.tex, source lines 1151--1151. Cited by the retained statement or proof but not part of the collected spans; supplied only as a navigation pointer. The argument does not depend on its content.\par\medskip
\ifdefined\hypertarget\hypertarget{sourcepointer-7}{}\fi
\noindent\textit{Source reference.} Original lemma 06HN: sources/stacks-artin/pinned-original-source.tex, source lines 1151--1151. Cited by the retained statement or proof but not part of the collected spans; supplied only as a navigation pointer. The argument does not depend on its content.\par\medskip
\ifdefined\hypertarget\hypertarget{sourcepointer-8}{}\fi
\noindent\textit{Source reference.} Original lemma 07WB: sources/stacks-artin/pinned-original-source.tex, source lines 1151--1151. Cited by the retained statement or proof but not part of the collected spans; supplied only as a navigation pointer. The argument does not depend on its content.\par\medskip
\ifdefined\hypertarget\hypertarget{sourcepointer-9}{}\fi
\noindent\textit{Source reference.} Original lemma 07WQ: sources/stacks-artin/pinned-original-source.tex, source lines 743--746. Cited by the retained statement or proof but not part of the collected spans; supplied only as a navigation pointer. The argument does not depend on its content.\par\medskip
\ifdefined\hypertarget\hypertarget{sourcepointer-10}{}\fi
\noindent\textit{Source reference.} Original lemma 07WX: sources/stacks-artin/pinned-original-source.tex, source lines 903--917. Cited by the retained statement or proof but not part of the collected spans; supplied only as a navigation pointer. The argument does not depend on its content.\par\medskip
\ifdefined\hypertarget\hypertarget{sourcepointer-11}{}\fi
\noindent\textit{Source reference.} Original lemma 07WU: sources/stacks-artin/pinned-original-source.tex, source lines 824--833. Cited by the retained statement or proof but not part of the collected spans; supplied only as a navigation pointer. The argument does not depend on its content.\par\medskip
\ifdefined\hypertarget\hypertarget{sourcepointer-12}{}\fi
\noindent\textit{Source reference.} Original lemma 07WL: sources/stacks-artin/pinned-original-source.tex, source lines 570--590. Cited by the retained statement or proof but not part of the collected spans; supplied only as a navigation pointer. The argument does not depend on its content.\par\medskip
\ifdefined\hypertarget\hypertarget{sourcepointer-13}{}\fi
\noindent\textit{Source reference.} Original lemma 07WK: sources/stacks-artin/pinned-original-source.tex, source lines 539--548. Cited by the retained statement or proof but not part of the collected spans; supplied only as a navigation pointer. The argument does not depend on its content.\par\medskip
\ifdefined\hypertarget\hypertarget{sourcepointer-14}{}\fi
\noindent\textit{Source reference.} Original lemma 07WV: sources/stacks-artin/pinned-original-source.tex, source lines 866--885. Cited by the retained statement or proof but not part of the collected spans; supplied only as a navigation pointer. The argument does not depend on its content.\par\medskip
\ifdefined\hypertarget\hypertarget{sourcepointer-15}{}\fi
\noindent\textit{Source reference.} Original lemma 07T2: sources/stacks-artin/pinned-original-source.tex, source lines 409--443. Cited by the retained statement or proof but not part of the collected spans; supplied only as a navigation pointer. The argument does not depend on its content.\par\medskip
\ifdefined\hypertarget\hypertarget{sourcepointer-16}{}\fi
\noindent\textit{Source reference.} Original lemma 07W4: sources/stacks-artin/pinned-original-source.tex, source lines 1151--1151. Cited by the retained statement or proof but not part of the collected spans; supplied only as a navigation pointer. The argument does not depend on its content.\par\medskip
\begin{lemma}
\label{lemma-finite-dimension}
Let $S$ be a locally Noetherian scheme. Assume
\begin{enumerate}
\item $\mathcal{X}$ is an algebraic stack,
\item $U$ is a scheme locally of finite type over $S$, and
\item $(\Sch/U)_{fppf} \to \mathcal{X}$ is a smooth surjective
morphism.
\end{enumerate}
Then, for any $\mathcal{F} = \mathcal{F}_{\mathcal{X}, k, x_0}$ as in
Section \ref{section-predeformation-categories}
the tangent space $T\mathcal{F}$ and infinitesimal automorphism space
$\text{Inf}(\mathcal{F})$ have finite dimension over $k$.
\end{lemma}

\begin{proof}
Let us write $\mathcal{U} = (\Sch/U)_{fppf}$. By our definition
of algebraic stacks the $1$-morphism $\mathcal{U} \to \mathcal{X}$
is representable by algebraic spaces. Hence in particular the
2-fibre product
$$
\mathcal{U}_{x_0} = (\Sch/\Spec(k))_{fppf} \times_\mathcal{X} \mathcal{U}
$$
is representable by an algebraic space $U_{x_0}$ over $\Spec(k)$. Then
$U_{x_0} \to \Spec(k)$ is smooth and surjective (in particular $U_{x_0}$
is nonempty). By Spaces over Fields, Lemma
\ref{spaces-over-fields-lemma-smooth-separable-closed-points-dense}
we can find a finite extension $l/k$ and a point
$\Spec(l) \to U_{x_0}$ over $k$. We have
$$
(\mathcal{F}_{\mathcal{X}, k , x_0})_{l/k} =
\mathcal{F}_{\mathcal{X}, l, x_{l, 0}}
$$
by Lemma \ref{lemma-change-of-field} and the fact that $\mathcal{X}$
satisfies (RS). Thus we see that
$$
T\mathcal{F} \otimes_k l \cong T\mathcal{F}_{\mathcal{X}, l, x_{l, 0}}
\quad\text{and}\quad
\text{Inf}(\mathcal{F}) \otimes_k l \cong
\text{Inf}(\mathcal{F}_{\mathcal{X}, l, x_{l, 0}})
$$
by
Formal Deformation Theory, Lemmas
\ref{formal-defos-lemma-tangent-space-change-of-field} and
\ref{formal-defos-lemma-inf-aut-change-of-field}
(these are applicable by
Lemmas \ref{lemma-algebraic-stack-RS} and
\ref{lemma-deformation-category} and
Remark \ref{remark-deformation-category-implies}).
Hence it suffices to prove that $T\mathcal{F}_{\mathcal{X}, l, x_{l, 0}}$
and $\text{Inf}(\mathcal{F}_{\mathcal{X}, l, x_{l, 0}})$
have finite dimension over $l$. Note that $x_{l, 0}$ comes from a point
$u_0$ of $\mathcal{U}$ over $l$.

\medskip\noindent
We interrupt the flow of the argument to show that the lemma for
infinitesimal automorphisms follows from the lemma for tangent spaces.
Namely, let
$\mathcal{R} = \mathcal{U} \times_\mathcal{X} \mathcal{U}$.
Let $r_0$ be the $l$-valued point $(u_0, u_0, \text{id}_{x_0})$ of
$\mathcal{R}$. Combining
Lemma \ref{lemma-fibre-product-deformation-categories} and
Formal Deformation Theory, Lemma
\ref{formal-defos-lemma-deformation-functor-diagonal}
we see that
$$
\text{Inf}(\mathcal{F}_{\mathcal{X}, l, x_{l, 0}})
\subset
T\mathcal{F}_{\mathcal{R}, l, r_0}
$$
Note that $\mathcal{R}$ is an algebraic stack, see
Algebraic Stacks, Lemma \ref{algebraic-lemma-2-fibre-product-general}.
Also, $\mathcal{R}$ is representable by an algebraic space $R$
smooth over $U$ (via either projection, see
Algebraic Stacks, Lemma \ref{algebraic-lemma-stack-presentation}).
Hence, choose an scheme $U'$ and a surjective \'etale morphism
$U' \to R$ we see that $U'$ is smooth over $U$, hence locally of
finite type over $S$. As $(\Sch/U')_{fppf} \to \mathcal{R}$ is
surjective and smooth, we have reduced the question to the case
of tangent spaces.

\medskip\noindent
The functor (\ref{equation-functoriality})
$$
\mathcal{F}_{\mathcal{U}, l, u_0}
\longrightarrow
\mathcal{F}_{\mathcal{X}, l, x_{l, 0}}
$$
is smooth by Lemma \ref{lemma-formally-smooth-on-deformation-categories}.
The induced map on tangent spaces
$$
T\mathcal{F}_{\mathcal{U}, l, u_0}
\longrightarrow
T\mathcal{F}_{\mathcal{X}, l, x_{l, 0}}
$$
is $l$-linear (by
Formal Deformation Theory, Lemma
\ref{formal-defos-lemma-k-linear-differential})
and surjective (as smooth maps of predeformation categories induce
surjective maps on tangent spaces by
Formal Deformation Theory, Lemma
\ref{formal-defos-lemma-smooth-morphism-essentially-surjective}).
Hence it suffices to prove that the tangent space of the deformation
space associated to the representable algebraic stack $\mathcal{U}$
at the point $u_0$ is finite dimensional. Let $\Spec(R) \subset U$ be
an affine open such that $u_0 : \Spec(l) \to U$ factors through $\Spec(R)$
and such that $\Spec(R) \to S$ factors through $\Spec(\Lambda) \subset S$.
Let $\mathfrak m_R \subset R$ be the kernel of the $\Lambda$-algebra map
$\varphi_0 : R \to l$ corresponding to $u_0$. Note that $R$, being of finite
type over the Noetherian ring $\Lambda$, is a Noetherian ring. Hence
$\mathfrak m_R = (f_1, \ldots, f_n)$ is a finitely generated ideal.
We have
$$
T\mathcal{F}_{\mathcal{U}, l, u_0}
=
\{\varphi : R \to l[\epsilon] \mid
\varphi \text{ is a } \Lambda\text{-algebra map and }
\varphi \bmod \epsilon = \varphi_0\}
$$
An element of the right hand side is determined by its values on
$f_1, \ldots, f_n$ hence the dimension is at most $n$ and we win.
Some details omitted.
\end{proof}

\end{document}
